\section{Real-data validation: battery SOC on held-out drive cycles}\label{sec:realdata}

C1 and C3 were tested once on real hardware with a design pre-registered before any training (Appendix~\ref{app:realdata}): the Panasonic 18650PF drive cycles \citep{kollmeyer2018panasonic} at 25$^\circ$C, split by cycle: Cycles 1--3, UDDS and LA92 for training, Cycle 4 for validation (model selection, rank-4-vs-dense-ablation choice, filter tuning), and US06, HWFET and NN as test cycles reported once. The target is the laboratory coulomb count ($Q=\RealQAh$\,Ah); the modeller's model is a constant-parameter 2-RC circuit identified by least squares on the training cycles (open-loop voltage RMSE \RealVmFitTrainMin--\RealVmFitTrainMax\,mV there, \RealVmFitValTestMin--\RealVmFitValTestMax\,mV on validation/test), used for phase~1 and inside the EKF/UKF ($Q,R$ tuned on the validation cycle). Episodes are 30-min windows through each discharge from a prior SOC 20\% off; the metric is settled SOC RMSE; a cycle counts only if the 95\% block-bootstrap CI of the paired difference excludes 0 and all \RealNSeeds\ seeds agree.

Both pre-registered rules are met (Table~\ref{tab:real}, Appendix~\ref{app:realdata}); they were decided on ``WS-best'', the per-seed validation choice between HyperKKL ($r=4$) and its dense ablation. H1 (C3): WS-best beats the tuned EKF on \RealHOneLearnedBetter\ of \RealNTestCycles\ test cycles (US06 \RealWsBestEkfUSZeroSixDiff, HWFET \RealWsBestEkfHWFETDiff, NN \RealWsBestEkfNNDiff\,\% SOC; likewise against the UKF), and HyperKKL at $r=4$ alone on \RealLoraEkfCyclesABetter\ of \RealNTestCycles. H2 (C1): WS-best beats latent injection on \RealHTwoWsBetter\ of \RealNTestCycles\ (by \RealWsBestLatentDiffMinAbs--\RealWsBestLatentDiffMaxAbs\,\% SOC, a small margin); HyperKKL at $r=4$ alone passes the strict rule on \RealLoraLatentCyclesABetter\ of \RealNTestCycles\ cycles with the same sign everywhere, and the simulated ordering weight-space $>$ latent $>$ concat $\gg$ autonomous holds. The mechanism is that of C3: on the aggressive US06 cycle the filters settle on a biased SOC (\RealEkfUSZeroSixRmse\,\% against \RealWsBestUSZeroSixRmse\,\%) because the constant-parameter circuit cannot explain the voltage error; the decoder trained on real data can. The limits matter as much: on whole discharges the filters have hours to converge and the gap closes except on US06 (EKF \RealWholeEkfNN\ vs learned \RealWholeWsBestNN\,\% on NN), so on mild cycles the advantage is mostly faster convergence; at 0$^\circ$C, where nothing was trained or identified, the Arrhenius-scaled EKF is better on \RealZeroEkfBetter\ of \RealZeroNCycles\ cycles (learned \RealZeroWsRmseMin--\RealZeroWsRmseMax\,\%); one cell, one campaign; a coulomb-counting reference; a quickly identified circuit model. The same-data baselines (Koopman KF, grey-box EKF) on the real cycles are \textcolor{red}{\textbf{\RealSameDataBaselines}}.
% TODO (pending): when experiments/real/results.json contains koopman_kf / greybox_* rows, re-run
% analysis/real_numbers.py (\RealSameDataBaselines flips to "present"), and report the comparison from
% Table tab:real; never claim the KKL observer is the best learner unless that table shows it.
